\documentclass[11pt,a4paper]{article}

% ---------- typography (identique au rapport principal) ----------
\usepackage[T1]{fontenc}
\usepackage[utf8]{inputenc}
\usepackage[french]{babel}
\usepackage{newtxtext,newtxmath}
\usepackage{microtype}

% ---------- layout ----------
\usepackage{geometry}
\geometry{a4paper, top=2.6cm, bottom=2.6cm, left=2.6cm, right=2.6cm}
\usepackage{xcolor}
\usepackage{booktabs}
\usepackage{tabularx}
\usepackage{enumitem}
\usepackage{needspace}

\widowpenalty=10000
\clubpenalty=10000

% ---------- hyperref (last) ----------
\usepackage[
    colorlinks=true,
    linkcolor=black,
    urlcolor={black!70},
    bookmarks=true,
    unicode=true
]{hyperref}
\usepackage{xurl}

\hypersetup{
    pdftitle={Legally Subjective : resume francais du projet de recherche},
    pdfauthor={Amine Harch El Korane, Ghali El Alj},
    pdfsubject={Resume en francais du projet Legally Subjective},
    pdfkeywords={Cour supreme, prediction, science des donnees, lycee, CGenial}
}

\begin{document}

% ================= bloc de titre (composé à la main, pas de \maketitle) =================
\begin{center}
{\LARGE\bfseries Legally Subjective\par}
\vspace{4pt}
{\large Jusqu'où peut-on prédire les décisions de la Cour suprême\\
des États-Unis à partir de données publiques ?\par}
\vspace{10pt}
{\normalsize Amine Harch El Korane\quad Ghali El Alj\par}
\vspace{2pt}
{\small Élèves de lycée (France) --- expériences, code et reproductibilité /
analyse juridique et lecture doctrinale.\par}
\vspace{2pt}
{\small Code et données : \href{https://github.com/Vitalcheffe/legally-subjective}{github.com/Vitalcheffe/legally-subjective}\par}
\vspace{4pt}
{\small Résumé en français du rapport de recherche complet (en anglais, 22 pages)\par}
\end{center}
\vspace{4pt}

\begin{center}
\rule{0.28\linewidth}{0.6pt}
\end{center}
\vspace{6pt}

\section*{La question, et pourquoi elle compte}

La Cour suprême des États-Unis tranche les questions les plus divisives de la
société américaine : avortement, port d'armes, discrimination, pouvoirs
présidentiels. Ses neuf juges sont nommés à vie précisément pour que leur
jugement ne soit \emph{pas} celui d'une machine : la subjectivité judiciaire
est, en théorie, le dernier rempart contre l'automatisation de la justice.
Nous nous sommes donc demandé jusqu'où cette subjectivité résiste à la
mesure : peut-on prédire la décision d'une affaire, ou le vote individuel
d'un juge, à partir d'informations publiques --- les votes passés, le texte
du dossier, le style rédactionnel du juge ? Les deux issues possibles
publient : si la prédiction marche, une part importante du jugement est
statistique ; si elle échoue sur les cas difficiles, alors la part
irréductible --- celle qui fait qu'un juge n'est pas un algorithme --- est
réellement là. Notre réponse est chiffrée des deux côtés, et c'est ce qui
fait la valeur du projet.

\section*{Les données : un corpus public, gelé et vérifiable}

Nous avons assemblé un corpus de \textbf{569 affaires plaidées}
(sessions OT2015 à OT2023), couvrant \textbf{13 juges} et
\textbf{4\,569 votes individuels codés}, à partir de deux sources publiques :
la \emph{Supreme Court Database} (votes et décisions) et
\emph{CourtListener} (textes des opinions). Le corpus a ensuite été
\emph{gelé} : chaque fichier est accompagné de son empreinte
cryptographique SHA-256, ce qui permet à quiconque de vérifier que les
données utilisées dans les résultats sont exactement celles qui ont été
publiées, et qu'aucune ligne n'a été modifiée après coup. L'intégralité du
projet a été menée à budget strictement nul : données publiques, matériel
de calcul standard et outils gratuits. Ce point n'est pas anecdotique : il
démontre qu'une démarche scientifique complète est accessible à des élèves,
sans laboratoire ni financement.

\section*{La méthode : une barre difficile, puis une épreuve scellée}

La démarche comporte deux temps. \emph{Premièrement}, fixer une barre
honnête : avant de tester un modèle de langage, nous avons mesuré ce que
donne une règle simple --- pour chaque juge, prédire systématiquement le
côté (conservateur ou libéral) que ce juge favorise d'habitude. Cette règle
idéologique atteint \textbf{63,7\,\% des votes individuels} sur la fenêtre
d'évaluation (intervalle de confiance à 95\,\% : [61,3 ; 65,9]). Tout
challenger doit battre ce score pour prétendre apporter quelque chose.
\emph{Deuxièmement}, l'épreuve scellée : nous avons mis de côté 50 décisions
serrées (résolues 5--4) mises sous scellés cryptographiques. Quatre
conditions ont produit leurs prédictions \emph{avant} que la moindre réponse
ne soit consultée, puis les prédictions elles-mêmes ont été hachées et
publiées, et la correction n'a été exécutée qu'une seule fois :
\begin{itemize}[leftmargin=1.4em,itemsep=2pt,topsep=4pt]
  \item \textbf{A --- zéro-shot} : un modèle de langage lit le dossier et
  prédit l'issue ;
  \item \textbf{B --- persona} : le même modèle, auquel on a « appris » un
  juge particulier, via un prompt biographique et un adaptateur affiné
  (QLoRA) sur les opinions passées de ce juge ;
  \item \textbf{C --- retrieval} : le dossier enrichi de précédents
  judiciaires récupérés automatiquement ;
  \item \textbf{D --- prior idéologique re-ajusté} : la règle simple,
  recalibrée sur les données les plus récentes.
\end{itemize}

\section*{Les résultats clés}

\begin{center}
\begin{tabular}{lcc}
\toprule
Condition & Décision de l'affaire & Vote individuel \\
\midrule
A --- zéro-shot & 61,22\,\% (30/49) & 56,23\,\% (203/361) \\
B --- persona (QLoRA) & 59,18\,\% (29/49) & 59,83\,\% (216/361) \\
C --- retrieval & 53,06\,\% (26/49) & 50,14\,\% (181/361) \\
D --- prior idéologique & 50,00\,\% (24/48) & \textbf{75,91\,\%} (312/411) \\
\bottomrule
\end{tabular}
\end{center}

Le message principal est un \emph{résultat négatif honnête} : sur les
cinquante affaires les plus serrées, aucune condition de langage ne bat la
barre idéologique, et la personnalisation ne détecte aucune amélioration
(persona 59,18\,\% contre zéro-shot 61,22\,\% ; test exact de McNemar,
$p=1{,}00$). Le prior idéologique, lui, s'effondre à pile-ou-face au
niveau de l'affaire (50\,\%) tout en dominant toutes les conditions de
langage au niveau du vote (75,9\,\%) : savoir \emph{qui} vote reste plus
informatif que savoir \emph{comment} le juge écrit. Les découvertes
secondaires sont tout aussi instructives : les neuf personae simulées
votent toutes plus conservatrices que le juge qu'elles imitent ; le prompt
biographique, et non l'adaptateur affiné, fait l'essentiel du
déplacement des votes ; l'hypothèse stylistique pré-enregistrée sur la
qualité des personae est contredite par les données (corrélation de
Spearman $\rho=-0{,}72$) ; et les personae se trompent sur les affaires
mêmes sur lesquelles elles ont été entraînées.

\section*{Ce que nous en concluons --- et nos limites assumées}

Une part mesurable de la décision judiciaire --- environ six votes sur dix
--- est extractible de l'historique public d'un juge ; mais sur les
affaires qui se jouent à une voix, cette statistique s'évanouit, et le
style rédactionnel d'un juge n'ajoute rien de détectable. Nous affichons
nos limites comme une composante du travail scientifique : avec 49 affaires
appariées, seule une amélioration de l'ordre de quinze à vingt points
aurait été détectable ; une partie des affaires scellées était potentiellement déjà
connue du modèle de base (contamination du pré-entraînement), même si le
contraste principal --- persona contre zéro-shot --- partage la même base
et annule cette contamination en différence ; le pipeline informatique a
été audité par un seul des deux auteurs ; et la chronologie fine de
l'analyse stylistique ne peut pas être établie par les seuls horodatages
git. Ces limites sont discutées une à une dans le rapport complet.

\section*{Reproductibilité et accès}

Tout est public et rejouable : code Python, corpus gelé, prédictions
scellées, journaux et scripts d'audit sur le dépôt GitHub du projet. Un
audit interne hostile a re-calculé chaque chiffre de ce résumé à partir
des artefacts primaires, avec concordance à quatre décimales près. La
science qui ne montre pas ses données n'est pas de la science ; nous avons
donc construit le projet pour que chaque affirmation de cette page puisse
être vérifiée, réfutée ou rejouée par un tiers, y compris avec les moyens
d'un lycée.

\begin{center}
\rule{0.28\linewidth}{0.6pt}
\end{center}
\vspace{4pt}

\begin{center}
{\small \emph{Deux élèves de lycée, un protocole scientifique complet :
question, données gelées, baselines pré-enregistrées, épreuve scellée,
statistiques appariées, audit indépendant --- et un résultat négatif
assumé.} \textbf{Legally Subjective}, github.com/Vitalcheffe/legally-subjective\par}
\end{center}

\end{document}
